\section{Experiment}
\subsection{Datasets}
We evaluate our method on two datasets:
GoEmotions~\cite{demszky-etal-2020-goemotions}, and Empathetic Dialogues~\cite{rashkin-etal-2019-towards}. 
Given our primary objective of fine-grained emotion classification, which requires distinguishing labels with subtle differences in meaning, we choose GoEmotions and Empathetic Dialogues for evaluation. These datasets are considered the most challenging datasets, as they contain a larger number of emotion labels with similar semantics.
Below we give the descriptions of these datasets.
GoEmotions is composed of comments from
Reddit~\cite{demszky-etal-2020-goemotions}. The total number of samples is 54k,
and each sample is annotated with one or multiple labels among 27 emotions and
neutral. To ensure a fair comparison with previous work~\cite{suresh-ong-2021-negatives}, 
we use only single-labeled samples and
exclude neutral ones. 
%Therefore, the remaining samples roughly account for 80\% of the whole dataset, and 
The training/validation/test split of the remaining dataset is 23,485\,/\,2,956\,/\,2,984.

Empathetic Dialogues~\cite{rashkin-etal-2019-towards} consists of conversations
with single emotion labels.  The situation was written by a speaker given an
emotion label. The listener was to respond with an empathetic utterance.
The process could go on for up to six turns.  Since the situation was written
based on the provided emotion, we used the situation as the model's input, 
following~\citet{suresh-ong-2021-negatives}. The dataset contains 24,850
conversations labeled among 32 emotions. The training/validation/test split of
the dataset is 19,533\,/\,2,770\,/\,2,547, respectively. Below, we use GE to
represent GoEmotions and ED to represent Empathetic Dialogues.

% Apart from applying the method to emotion dataset, we also incorporated two news datasets, The New York Times (NYT) and The 20 Newsgroups (20News) in our experiment. The label set in these two datasets also has hierarchical structure. Therefore, we can see that if our method can generalize to another domain. The descriptions of these two datasets are provided as follows.

% The New York Times (NYT) contains articles published on The New York Times. Each article can be classified into 5 coarse-grained labels (e.g., politics, science) and 26 fine-grained labels (e.g., gun control, cosmos). 

% The 20 Newsgroups (20News)\footnote{\url{http://qwone.com/~jason/20Newsgroups/}} is a collection of news articles. Each article can be  categorized into one of 6 coarse-grained genres (e.g., science, computers) and 17 fine-grained labels (e.g., medicine, space, graphics).
% Statistics of the above-mentioned datasets are listed on Table~\ref{tab:datasets}.

% \begin{table}[t]
% \centering
% \begin{tabular}{cll}
% \bottomrule
% Dataset          & \multicolumn{1}{c}{$|\mathcal{D}|$} & \multicolumn{1}{c}{$|\mathcal{M}|$} \\ \hline
% GoEmotions       & 29,425                  & 27                      \\
% EmpatheticDialog & 24,850                  & 32                      \\
% NYT              & 11,744                  & 26                      \\
% 20News           & 16,468                  & 17\\
% \toprule
% \end{tabular}

% \caption{Statistics of the datasets.}
% \label{tab:datasets}
% \end{table}

\subsection{Experiment Settings and Baselines}
We compare the proposed \wsmodelname primarily with three categories of strong
baselines:

\noindent\textbf{General pre-trained language models.} 
We compared against BERT~\cite{devlin-etal-2019-bert}, RoBERTa~\cite{roberta},
and ELECTRA~\cite{clark2020electra}. These models are pre-trained on large unlabeled corpora. They generate high-quality representations, and all perform well on text classification tasks. We also compare different size of these pre-trained models, denoted as $\mathrm{base}$ and $\mathrm{large}$.

\noindent\textbf{Label Embedding-aware models.}
\citet{suresh-ong-2021-negatives} propose label-aware contrastive loss (LCL),
which weights each negative sample differently. Specifically, more confusable pairs contribute more to the objective function; this yields promising results on fine-grained text classification.
In addition, we compare against HiAGM~\cite{jie2020hierarchy}, the strongest hierarchy-aware text classification model with source code publicly available.
Lastly, we implement a baseline called LabelEmb, which encodes the label description (i.e., the definition of emotions) to derive label embeddings, and train the model on the Euclidean space with the rest setting same as \modelname.

\noindent\textbf{Hyperbolic classification models.}
We also compared with models trained in hyperbolic space for classification, including (1)~Hyperbolic SVM (HSVM) proposed by~\citet{cho2019large}, which
generalizes the support vector machine to hyperbolic space, and (2)~Hyperbolic
Neural Model (HNN) proposed by~\citet{ganea2018hyperbolic}, a hyperbolic GRU that performs all necessary operations in hyperbolic space to train a neural
network. (3)~Hyperbolic Interaction Model (HyperIM)~\cite{chen2020hyperbolic} jointly learns word and label embeddings, and measure the similarities in the Poincor\'{e} disk to aggregate input representations.
(4)~\textsc{HiddeN}~\cite{chatterjee-etal-2021-joint} is a framework which does not assume the label hierarchy is known. It also proposed to learn the label embedding jointly in an end-to-end fashion.
% (5)~HyperHAN~\cite{ijcai2020p552}
% (6)~HyperText~\cite{zhu-etal-2020-hypertext}
For a fair comparison, we set the word dimension of the hyperbolic 
space to 100, the same as 
the dimension we use.  % AMH: what is this?
% There are a few works trying to address sentiment analysis~\cite{DBLP:journals/corr/abs-2012-01786, yin-etal-2020-sentibert, ke-etal-2020-sentilare} and sentiment recognition~\cite{suresh-ong-2021-negatives}. Therefore, we attempted to conduct experiments on some previous works and compare the results with \modelname. 

\noindent\textbf{Evaluation metrics.}
Following~\citet{suresh-ong-2021-negatives}, we use accuracy and weighted F1 as
the evaluation metrics. Weighted F1 takes into account the number of samples in each
class, and weights the macro F1 by this ratio. This can be expressed as
\begin{equation}
    F1_{\mathit{weighted}} = 2 \sum_c \frac{n_c}{N} \frac{P_c \times R_c}{P_c + R_c},
\end{equation}
where $n_c$ is the number of samples in class~$c$, $N$ is the number of total
samples, and $P_c$ and $R_c$ are the precision and recall for class~$c$, respectively.
\begin{comment}
	\begin{table}[t]
		 \centering
		 \scalebox{0.88}{
		 \begin{tabular}{|c c c c c c|}
		 \hline
				\multicolumn{2}{|c}{\small Datasets} & \multicolumn{2}{c}{\small ED} & \multicolumn{2}{c|}{\small GE}  \\
				\multicolumn{2}{|c}{\small Number of classes} & \multicolumn{2}{c}{\small 32} & \multicolumn{2}{c|}{\small 27} \\
				\hline
				\multicolumn{2}{|c}{} & \textbf{\small Acc} & \textbf{\small F1} & \textbf{\small Acc} & \textbf{\small F1} \\
				\hline
				\multicolumn{2}{|c}{\small SentiLARE ~\cite{ke-etal-2020-sentilare}} & \small 58.6 & \small 59.4 & \small \textcolor{red}{65.6}  & \small \textcolor{red}{66.3}    \\
				\hline
				\multicolumn{2}{|c}{\small SentiBERT ~\cite{yin-etal-2020-sentibert}} &  & & &   \\
				\hline
				\multicolumn{2}{|c}{\small Self-Explaining ~\cite{DBLP:journals/corr/abs-2012-01786}} &  \small \textcolor{red}{57.9} & \small 58.4  & \small 64.3  & \small 65.2 \\
				\hline
				\multicolumn{2}{|c}{\small LCL ~\cite{suresh-ong-2021-negatives}} & \small \textcolor{red}{60.1}  &  \small 59.1 &  \small\textcolor{red}{65.5} & \small 64.8  \\
				\hline
				\multicolumn{2}{|c}{\small HypoEmo} & \textbf{\small 57.9} & \textbf{\small 60.6} & \textbf{\small 65.6}  & \textbf{\small 66.5} \\
				\hline
		 \end{tabular}}
		 \caption{Accuracy and weighted f1 score evaluated on a test set of ED and GE among three different baselines and our proposed method \wsmodelname}
		 \label{tab:baseline}
	\end{table}
\end{comment}

\noindent\textbf{Implementation Details.} \wsmodelname is a encoder-agnostic framework which can easily adapt to different kinds of text encoders. 
In the experiment, we use pre-trained $\mathrm{RoBERTa}_\mathrm{base}$ as the backbone, which has 12 layers with a hidden size of 768. 
%Due to the limited computational power, we adopt $\mathrm{RoBERTa}_\mathrm{base}$, 
%In addition, to project text representation extracted from the model to the hyperbolic space, we used Hyperbolic Neural Network ~\cite{ganea2018hyperbolic} to retrieve the text embedding in the hyperbolic space. Therefore, we could compute the distance between the text embedding and the pretrained label embedding in the hyperbolic space.
%In the augmentation step, we used back-translation ~\cite{sennrich-etal-2016-improving}, and the augmented sample is in  the positive set of the original sample.
During training, we applied the Adam optimizer in Euclidean space
with a learning rate of $10^{-5}$ and a weight decay of $0.01$. 
By contrast, we utilized Riemannian Adam~\cite{becigneul2018riemannian_adam} 
to train our label
embeddings in hyperbolic space with a learning rate of $0.01$. The dimension of hyperbolic label embedding is set to 100, which is searched from \{2, 10, 50, 100, 250\}. Other implementation details can be found in our code.

\subsection{Main Results}
\label{sec:main_results}
\noindent\textbf{Baseline comparison.} To demonstrate the effectiveness of \modelname, we conduct experiments to compare the performance of different models. The comparison is shown in Table~\ref{tab:main_results}. Firstly, we compare \wsmodelname with general pre-trained language models. Among them, $\mathrm{RoBERTa}_\mathrm{large}$ performs the best, while \wsmodelname outperforms it on weighted F1 by 2.8\% on ED and 1.1\% on GE. It is also worth mentioning that \wsmodelname has considerably smaller parameter size compared with  $\mathrm{RoBERTa}_\mathrm{large}$ (125M v.s. 355M), resulting in significantly lower training and inference time.
This indicates the effectiveness of the proposed label-aware hyperbolic embeddings and the strategy to weight the standard cross entropy loss by hyperbolic distance.  

Secondly, we compare with label-aware system. Since LCL is the previous state-of-the-art, we mainly compare the efficiency with it ($\Delta$ in the left of Table~\ref{tab:main_results}). Results show that our proposed method outperforms LCL with much higher efficiency. This is because LCL augments data by using the synonym replacement technique, which doubles the size of data. Also, they use two encoders to train the main classifier and a weighting network, which doubles the parameter size. In contrast, \wsmodelname uses single encoder and uses only the original samples without any data augmentation method, and still out-wins LCL by 2.8\% and 2.5\% absolute F1 score on ED and GE, respectively. 
Moreover, Although HiAGM take into account the label hierarchy, it utilize RNN architecture, making it less efficient and underperforming \wsmodelname by a large margin.
Lastly, \wsmodelname performs better than LabelEmb, which LabelEmb calculates the weighted loss in the Euclidean space. This again demonstrates the efficacy of our proposed hyperbolic space integration. 

Also, we notice that \wsmodelname works better than models that are fully trained on hyperbolic space, which indicates the benefits of jointly learning the hyperbolic label embedding and fine-tuning $\mathrm{RoBERTa}_\mathrm{base}$ in a hybrid space settings. This hybrid setting could benefit from both the power of pre-trained language model and the strength of hyperbolic space to capture hierarchical information. To sum up, we could achieve better results compared to previous works without increasing data size or model parameters, which is more effective and efficient.  

\begin{table*}[t]
    \small
    \centering
    \begin{tabular}{c c c c c | c c c c }
    \bottomrule
    \\[-1em]
         \multicolumn{5}{c}{} & \multicolumn{2}{c}{ Empathetic Dialogues} & \multicolumn{2}{c}{ GoEmotions} \\ \\[-1em] \hline \\[-1em]
        %  \multicolumn{2}{c}{ Number of classes} & \multicolumn{2}{c|}{32} & \multicolumn{2}{c}{ 27} \\
         \multicolumn{1}{c}{} & \textbf{\#Params (\textdownarrow)} & $\xi_{tr}$ (\textdownarrow) & $\xi_{conv}$ (\textdownarrow) & $\xi_{inf}$ (\textdownarrow)& \textbf{ Acc} (\textuparrow)& \multicolumn{1}{c}{\textbf{ F1} (\textuparrow)}  & \textbf{ Acc} (\textuparrow) & \multicolumn{1}{c}{\textbf{F1} (\textuparrow)} \\ \\[-1em] \toprule 
        %  \multicolumn{2}{c}{$RoBERTA +\alpha L_{ce}$}  &  54.5 \footnotesize{$\pm$ 0.7}  &  \multicolumn{1}{c|}{56.0 \footnotesize{$\pm$ 0.4}}  &  62.6 \footnotesize{$\pm$ 0.6} &  \multicolumn{1}{c}{64.0 \footnotesize{$\pm$ 0.2}}  \\
        %  \hline
        %  \multicolumn{2}{c}{$RoBERTA + \alpha L_{wce}$} &  58.5 \footnotesize{$\pm$ 0.5}  & \multicolumn{1}{c|}{60.5 \footnotesize{$\pm$ 0.4}}  &  65.1 \footnotesize{$\pm$ 0.3}  & \multicolumn{1}{c}{65.8 \footnotesize{$\pm$ 0.3}}  \\
      
        %  \multicolumn{2}{c}{ $RoBERTA + \alpha L_{wce}+ \beta L_{pc}$} &  58.1 \footnotesize{$\pm$ 0.3}  &  \multicolumn{1}{c|}{60.1 \footnotesize{$\pm$ 0.3}}  & 65.1 \footnotesize{$\pm$ 0.2} & \multicolumn{1}{c}{65.9 \footnotesize{$\pm$ 0.2}} \\
         
         %\multicolumn{2}{c}{$RoBERTA + \alpha L_{wce}+ \beta L_{pc} + \gamma L_{cl}$} & 58.4 \footnotesize{$\pm$ 0.3} & \multicolumn{1}{c|}{60.3 \footnotesize{$\pm$ 0.3}} &  65.2 \footnotesize{$\pm$ 0.2} &  \multicolumn{1}{c}{66.0 \footnotesize{$\pm$ 0.2}} \\
         \multicolumn{1}{c}{$\mathrm{BERT}_\mathrm{base}$} & 110M & 96.7 & 290.0 & 1.1 & 50.4 \footnotesize{$\pm$0.3} & \multicolumn{1}{c}{51.8 \footnotesize{$\pm$0.1}} & 60.9 \footnotesize{$\pm$0.4} & \multicolumn{1}{c}{62.9 \footnotesize{$\pm$0.5}} \\

         \multicolumn{1}{c}{$\mathrm{RoBERTa}_\mathrm{base}$} & 125M & 99.3 & 297.9 & 1.2 & 54.5 \footnotesize{$\pm$0.7}  &  \multicolumn{1}{c}{56.0 \footnotesize{$\pm$0.4}}  &  62.6 \footnotesize{$\pm$0.6} &  \multicolumn{1}{c}{64.0 \footnotesize{$\pm$0.2}}\\
         
         \multicolumn{1}{c}{$\mathrm{ELECTRA}_\mathrm{base}$} & 110M & 97.6 & 292.7 & 1.1 & 47.7 \footnotesize{$\pm$1.2} & \multicolumn{1}{c}{49.6 \footnotesize{$\pm$1.0}} & 59.5 \footnotesize{$\pm$0.4} & \multicolumn{1}{c}{61.6 \footnotesize{$\pm$0.6}} \\
         
         \multicolumn{1}{c}{$\mathrm{BERT}_\mathrm{large}$} & 340M & 181.0 & 362.0 & 3.5 & 53.8 \footnotesize{$\pm$0.1} & \multicolumn{1}{c}{54.3 \footnotesize{$\pm$0.1}} & 64.5 \footnotesize{$\pm$0.3} & \multicolumn{1}{c}{\underline{65.2} \footnotesize{$\pm$0.4}} \\

         \multicolumn{1}{c}{$\mathrm{RoBERTa}_\mathrm{large}$} & 355M & 185.7 & 371.4 & 3.7 & 57.4 \footnotesize{$\pm$0.5}  &  \multicolumn{1}{c}{\underline{58.2} \footnotesize{$\pm$0.3}}  &  \underline{64.6} \footnotesize{$\pm$0.3} &  \multicolumn{1}{c}{\underline{65.2} \footnotesize{$\pm$0.2}}\\
         
         \multicolumn{1}{c}{$\mathrm{ELECTRA}_\mathrm{large}$} & 335M & 179.9 & 539.8 & 3.5 & 56.7 \footnotesize{$\pm$0.6} & \multicolumn{1}{c}{57.6 \footnotesize{$\pm$0.6}} & 63.5 \footnotesize{$\pm$0.3} & \multicolumn{1}{c}{64.1 \footnotesize{$\pm$0.4}} \\
         \hline
         % \multicolumn{3}{c}{SentiLARE~\cite{ke-etal-2020-sentilare}} & 125M & 59.2 \footnotesize{(0.4)} & \multicolumn{1}{c}{59.8 \footnotesize{(0.5)}} & \textbf{65.5} \footnotesize{(0.2)} & \multicolumn{1}{c}{\textbf{66.2} \footnotesize{(0.3)}}\\
                  
         % \multicolumn{3}{c}{Self-Explaining~\cite{DBLP:journals/corr/abs-2012-01786}} & 125M & 57.4 \footnotesize{(0.2)} & \multicolumn{1}{c}{57.6 \footnotesize{(0.2)}} & 63.4 \footnotesize{(0.5)} &  \multicolumn{1}{c}{63.7 \footnotesize{(0.5)}}\\
         
         % \multicolumn{3}{c}{ LCL~\cite{suresh-ong-2021-negatives}} & 220M &  \textbf{60.1} \footnotesize{(0.3)} & \multicolumn{1}{c}{ 59.1 \footnotesize{(0.3)}} &  \textbf{65.5} \footnotesize{(0.2)} & \multicolumn{1}{c}{ 64.8 \footnotesize{(0.2)}}\\
         \\[-8pt]
         \multicolumn{1}{c}{$LCL^\dag$} & 220M & 421.7 & 1264.9 & 3.9 & \underline{59.1} \footnotesize{$\pm$0.4} & \multicolumn{1}{c}{\underline{58.2} \footnotesize{$\pm$0.5}} & \underline{64.6} \footnotesize{$\pm$0.2} & \multicolumn{1}{c}{63.8 \footnotesize{$\pm$0.3}}\\
         
         \multicolumn{1}{c}{HiAGM} & 15M & 2673.2 & 10692.9 & 8.6 & 47.8 \footnotesize{$\pm$0.6} & \multicolumn{1}{c}{50.2 \footnotesize{$\pm$0.7}} & 59.7 \footnotesize{$\pm$0.6} & \multicolumn{1}{c}{61.8 \footnotesize{$\pm$0.5}}\\
         
         \multicolumn{1}{c}{LabelEmb} & 125M & 103.6 & 518.1 & 1.1 & 55.1 \footnotesize{$\pm$0.7} & \multicolumn{1}{c}{56.2 \footnotesize{$\pm$0.5}} & 62.7 \footnotesize{$\pm$0.6} & \multicolumn{1}{c}{62.8 \footnotesize{$\pm$0.4}}\\
         \hline
         % \multicolumn{3}{c}{ LCL (Proposed implementation, w/o aug)} & 220M & 59.6 \footnotesize{(0.4)} & \multicolumn{1}{c}{58.7 \footnotesize{(0.4)}} & 64.7 \footnotesize{(0.5)} & \multicolumn{1}{c}{64.0 \footnotesize{(0.4)}}\\
         % \hline
         \\[-8pt]
         \multicolumn{1}{c}{HSVM} & 428K & 14.7 & 42.1 & 0.7 & 27.4 \footnotesize{$\pm$0.0} & \multicolumn{1}{c}{26.7 \footnotesize{$\pm$0.0}} & 23.6 \footnotesize{$\pm$0.0} & \multicolumn{1}{c}{22.3 \footnotesize{$\pm$0.0}}\\
         
         \multicolumn{1}{c}{HNN} & 5M & 15421.5 & 92529.1 & 17.5 & 41.2 \footnotesize{$\pm$0.9} & \multicolumn{1}{c}{42.0 \footnotesize{$\pm$0.8}} & 46.6 \footnotesize{$\pm$0.6} & \multicolumn{1}{c}{47.2 \footnotesize{$\pm$0.5}}\\

          \multicolumn{1}{c}{HyperIM} & 5M & 2266.3 & 6798.9 & 8.4 & 44.1 \footnotesize{$\pm$1.2} & \multicolumn{1}{c}{43.6 \footnotesize{$\pm$1.0}} & 50.2 \footnotesize{$\pm$0.9} & \multicolumn{1}{c}{49.7 \footnotesize{$\pm$0.7}}\\
          
          \multicolumn{1}{c}{\textsc{HiddeN}} & 11M & 13473.0 & 67364.8 & 16.6 & 42.9 \footnotesize{$\pm$1.4} & \multicolumn{1}{c}{44.3 \footnotesize{$\pm$1.1}} & 47.2 \footnotesize{$\pm$1.1} & \multicolumn{1}{c}{49.3 \footnotesize{$\pm$0.9}}\\
          
         \hline
         \\[-8pt]
          \multicolumn{1}{c}{\wsmodelname} & 125M & 97.6 & 585.8 & 1.2 & \textbf{59.6} \footnotesize{$\pm$0.3} & \multicolumn{1}{c}{\textbf{61.0} \footnotesize{$\pm$0.3}}  & \textbf{65.4} \footnotesize{$\pm$0.2}& \multicolumn{1}{c}{\textbf{66.3} \footnotesize{$\pm$0.2}} \\
          
          \multicolumn{1}{c}{$\Delta$} & -43.2\% & -76.9\% & -53.3\% & -68.5\% & +0.8\% & +4.8\% & \multicolumn{1}{c}{+0.8\%} & \multicolumn{1}{c}{+3.9\%}  \\
         \toprule

    \end{tabular}
	 \caption{Left: efficiency comparison (lower is better), where $\xi_{tr}$ denotes average training time for one epoch, $\xi_{conv}$ denotes the average time to achieve the best results, and $\xi_{inf}$ denotes the average inference time for the whole testing set. We evaluate efficiency for all systems on GE with batch size equals 16. Right: performance comparison (higher is better.) We conducted the experiment five times with different seeds, and report the average score with the standard deviation. The best (second-best) results are set in \textbf{bold} (\underline{underlined}). ``\dag'' denotes reproduced results from official implementation, and $\Delta$ represents the efficiency and performance gain compared with LCL, the previous state-of-the-art, in relative percentage.}
    \label{tab:main_results}
\end{table*}

\begin{table}[ht!]
\centering
\small
\begin{tabular}{clcc}
\bottomrule\\[-1em]
\multicolumn{1}{c}{Dataset} & \multicolumn{1}{c}{Model}   & \multicolumn{1}{c}{F1} & \multicolumn{1}{c}{F1 (+HypEmo)} \\\\[-1em] \toprule
\multirow{3}{*}{ED}  & $\mathrm{BERT}_\mathrm{base}$    & 51.8                   & \textbf{57.7}                    \\
                     & $\mathrm{RoBERTa}_\mathrm{base}$ & 56.0                   & \textbf{61.0}                    \\
                     & $\mathrm{ELECTRA}_\mathrm{base}$ & 57.6                   & \textbf{58.9}                    \\ \hline \\[-1em]
\multirow{3}{*}{GE}  & $\mathrm{BERT}_\mathrm{base}$    & 62.9                   & \textbf{65.3}                    \\
                     & $\mathrm{RoBERTa}_\mathrm{base}$ & 64.0                   & \textbf{66.3}                    \\
                     & $\mathrm{ELECTRA}_\mathrm{base}$ & 64.1                   & \textbf{65.7}                    \\ \toprule
\end{tabular}
\caption{Performance in terms of weighted F1 score when \wsmodelname is added on different encoders.}
\label{tab:encoder_analysis}
\end{table}

% \begin{table}[ht!]
% \small
% \centering
% \begin{tabular}{lcc}
% \bottomrule \\[-1em]
% Model            & Empathetic Dialogues & GoEmotions \\ \\[-1em] \toprule
% $\mathrm{BERT}_\mathrm{base}$             & 51.8 & 62.9 \\
%  + HypEmo    & \textbf{57.7} & \textbf{65.3} \\
% $\mathrm{RoBERTa}_\mathrm{base}$          & 56.0 & 64.0 \\
%  + HypEmo & \textbf{61.0} & \textbf{66.3} \\
% $\mathrm{ELECTRA}_\mathrm{base}$          & 57.6 & 64.1 \\
%  + HypEmo & \textbf{58.9} & \textbf{65.7} \\ \toprule
% \end{tabular}
% \caption{Performance in terms of weighted F1 score when \wsmodelname is added on different encoders.}
% \label{tab:encoder_analysis}
% \end{table}

\noindent\textbf{Performance on different encoder.}
We apply \wsmodelname on top of different encoders to examine whether \wsmodelname is a model-agnostic method that could bring improvement regardless of the encoder being used. 
Table~\ref{tab:encoder_analysis} shows the results in terms of weighted F1 score. We observe that no matter which encoder is adopted, adding \wsmodelname leads to further improvements. For instance, applying \wsmodelname on $\mathrm{BERT}_\mathrm{base}$ enhances the performance by 5.9\% absolute percentage on ED, and the same phenomenon can be observed on $\mathrm{RoBERTa}_\mathrm{base}$ and $\mathrm{ELECTRA}_\mathrm{base}$ across two datasets. This verifies that \wsmodelname is model-agnostic and could be easily built on top of any text encoder to boost performance.


\subsection{Case Study}
\begin{figure}[ht!]
    \centering
    \includegraphics[width=0.475\textwidth]{figures/ED_subset.png}
    \caption{Case study for various ED subsets. We report F1 for brevity. CE is the standard cross entropy loss, and SupCon represents supervised contrastive loss~\cite{NEURIPS2020_supcon}.}
    \label{fig:ED_subset}
\end{figure}
Following~\citet{suresh-ong-2021-negatives}, we investigate the performance of
\wsmodelname when the label set contains pairs with close semantics. We compare
the proposed \wsmodelname with different objectives and encoders, and follow
the ED subsets selected by \citet{suresh-ong-2021-negatives}, which
include the most difficult sets chosen after enumerating all combinations that contain
four labels. These subsets are a: \{\emph{Anxious, Apprehensive, Afraid,
Terrified}\}, b: \{\emph{Devastated, Nostalgic, Sad, Sentimental}\}, c:
\{\emph{Angry, Ashamed, Furious, Guilty}\}, and d: \{\emph{Anticipating, Excited,
Hopeful, Guilty}\} from the ED dataset. We conducted the experiments with
RoBERTA and ELECTRA with standard cross entropy loss and
supervised contrastive loss~\cite{NEURIPS2020_supcon}. 
For supervised contrastive loss, we adopt back translation~\cite{sennrich-etal-2016-improving} to form the positive pairs,
and we view it as a strong competitor because we expect contrastive loss helps to learn better representations could hence improve the FEC task.
The result is shown in Fig~\ref{fig:ED_subset}. First, \wsmodelname outperforms all baselines by a large margin on the most difficult subset (a), which demonstrates the advantage of incorporating hyperbolic space when addressing fine-grained classification task. In particular, \wsmodelname beats the previous state-of-the-art, LCL, on this subset with a substantial improvement (54.9 v.s. 63.6.) Also, \wsmodelname surpasses models with cross entropy loss and supervised contrastive loss on a / b / c and comparably on d. As a / b / c are the top three most difficult sets, this result shows that label-aware hyperbolic weighted loss is conducive to separation under label sets which are more confusing. 
In addition, we compare with LCL with and without augmentation to ensure a fair comparison. The result shows
that \wsmodelname consistently outperforms LCL even with data augmentation that doubles the size of samples. 
%This demonstrates the effectiveness of incorporating hyperbolic embeddings and the joint learning architecture to distinguish between semantically close pairs. 
In summary, the proposed model performs the best when the label set contains confusing classes. For the most challenging sets, \wsmodelname outperforms models trained with conventional cross entropy loss and supervised contrastive loss, and even the state-of-the-art, LCL, while being more sample-efficient. When the label set is simpler, \wsmodelname performs on par with the others.

%Although these subsets are the top four most challenging subsets, we expect proposed \wsmodelname can perform even better when the label set itself is more confusing. As shown in Table \ref{tab:ED_exp}, \wsmodelname performs comparably under 4-easy set, yet we can see the consistent improvement on all the hard sets. Specifically, we implemented two ~\citealt{suresh-ong-2021-negatives} models with and without augmentation. To make a fair comparison, we compared \wsmodelname with LCL without augmentation and the result showed that our model's performance was better than that of LCL. In addition, our model also out-wins LCL with augmentation. This exhibits the effectiveness of incorporating hyperbolic embeddings and the jointly learning architecture to distinguish between semantically close pairs. 
%Although supervised contrastive learning is known for learning well separated embeddings, it is more time-consuming and less efficient since it require positive pair for each sample and view all the other in the same batch as negative. Conversely, adopting hyperbolic embeddings is more efficient, and could perform even better than contrastive learning method. 

% \begin{table}[h]
% \small
% \centering
% \begin{tabular}{cccccc}
% \bottomrule \\[-1em]
% ED subset               & a & b & c & d \\ \\[-1em]
% \toprule
% RoBERTa + CE & 59.8  & 67.4 & 57.0  & 79.5 \\
% ELECTRA + CE & 57.3 & 68.0 & 57.8 & \textbf{80.7} \\
% RoBERTa + SupCon   & \underline{59.9} & \underline{69.1} & \textbf{60.3} & 79.9 \\
% ELECTRA + SupCon        & 57.7 & 66.8 & \underline{58.7} & \underline{80.1} \\
% LCL w/ aug       & 55.0 &  59.3 & 54.7 &  79.1 \\
% LCL w/o aug & 54.9 & 62.8 & 51.3 & 78.7  \\ \\[-1em]
% \hline \\[-1em]
% HypEmo                  & \textbf{63.6} & \textbf{69.5} & \textbf{60.3} & \underline{80.1} \\ 
% \toprule
% \end{tabular}
% 	 \caption{Case study for various ED subsets. We report only
% 	 F1 for brevity. CE is the standard cross entropy loss, and SupCon represents supervised contrastive
% 	 loss~\cite{NEURIPS2020_supcon}.}
%     \label{tab:ED_exp}
% \end{table}

\begin{figure*}[ht!]
\centering
\begin{subfigure}{.3\textwidth}
    \centering
    \includegraphics[width=4.75cm]{figures/tsne_ED_CE.png}
    \caption{Standard cross entropy loss}
    \label{fig:tsne1}
\end{subfigure}
\begin{subfigure}{.3\textwidth}
    \centering
    \includegraphics[width=4.75cm]{figures/tsne_ED_SupCon.png}
    \caption{Supervised contrastive loss}
    \label{fig:tsne2}
\end{subfigure}
\begin{subfigure}{.3\textwidth}
    \centering
    \includegraphics[width=4.75cm]{figures/tsne_ED_WCE.png}
    \caption{Hyperbolic weighted loss}
    \label{fig:tsne3}
\end{subfigure}
\caption{Representation generated from text encoder using different
objective functions. We take \emph{afraid} and \emph{terrified} as examples for semantically close labels, and we highlight their position with the yellow arrow. Evidently, our proposed label-aware hyperbolic weighted loss leads to larger
inter-class distances.}
\label{fig:visual}
\end{figure*}

\begin{figure}[]
\centering
\includegraphics[width=0.3\textwidth]{figures/hierarchy.png}
\caption{Label hierarchy method vs.\ performance.}
\label{fig:hierarchy}
\end{figure}

\subsection{Effect of Label Hierarchy}
In this section, we investigate how the method used
to form the label hierarchy affects performance. We compare three settings: \emph{None}, \emph{Random} and \emph{Parrott}, as shown in Fig.~\ref{fig:hierarchy}. 
\textit{None} means that we abandon all hierarchy and distribute labels
uniformly on the Poincar\'{e} ball, and \textit{Random} means that
we randomly shuffle the correct order to make hierarchy meaningless. 
\textit{Parrott} is the emotion model proposed by~\citet{parrott2001emotions},
which defines a three-level emotion hierarchy. We observe that following 
an expert-defined hierarchy yields the best performance, better than without hierarchical
information, and the worst is shuffling to discard all hierarchical information.
Accordingly, we follow Parrott's emotion model to form the label hierarchy and learn the hyperbolic label embeddings.

\subsection{Visualization of Representations}
To better understand the effect of the proposed label-aware hyperbolic weighted
loss, we compare representations learned with different commonly used
objective functions. We trained on the ED dataset and projected
the embeddings generated from the text encoder onto 2-dimensional space with
t-SNE~\cite{van2008visualizing}. We compare standard cross entropy loss
(Fig.~\ref{fig:tsne1}), supervised contrastive loss~\cite{NEURIPS2020_supcon}
(Fig.~\ref{fig:tsne2}), and the proposed label-aware hyperbolic weighted loss
(Fig.~\ref{fig:tsne3}). For supervised contrastive loss, we also used
back translation as the augmentation to form
positive pairs. Then, standard cross entropy loss was summed with supervised
contrastive loss. As data augmentation doubles the data size, we
expect the supervision and the objective to lead to more a separated
representation.
In Fig.~\ref{fig:tsne1}, we observe indistinct clustering, as
many data points with different classes are mingled together, showing the
difficulties of fine-grained classification in which different classes
often share close semantics. In particular, \emph{afraid} and \emph{terrified} are confusing for models and their representations are mixed. 
With the use of supervised contrastive loss shown
in Fig.~\ref{fig:tsne2}, the clustering becomes somewhat clearer but at the
cost of increasing data.
Last, in Fig.~\ref{fig:tsne3}, the inter-class distance is
clearly larger than others, and the clusters are also more dispersed. 
Specifically, the representations of \emph{afraid} and \emph{terrified} are much more separated.
This shows the advantage of the proposed label-aware hyperbolic loss, which yields
better representations, even without the need for additional significant costs.